\section{Inside the group: size and leader}\label{sec:inside}

Does the effect grow with the number of pilots, and does the pilot ahead fly differently from the one
behind?

\subsection{Size of the configuration}\label{sec:size}

The group time is cut where the configuration changes, so every stretch has a fixed number of pilots
$n$ (2, 3, 4, 5--6, 7+), and a flight can contribute to several sizes. Two contrasts of
Eq.~\eqref{eq:contrast} follow the size: $\G_n-\A$ within the flight, and $\G_n-\G_2$ at the same
site-day and hour, which asks whether a group does more than a pair.

\begin{figure}
\centering
\includegraphics[width=0.56\textwidth]{size_hloc.pdf}
\caption{$\Hloc(\tau)$ over every lag by the size of the configuration (orange, darker with more
pilots), $\Sl$ (blue) and, dashed, the same flights alone for each size.}
\label{fig:size-hloc}
\end{figure}

\begin{table}
\centering\footnotesize\setlength{\tabcolsep}{3pt}
\begin{tabular}{@{}l r c c c c c c c c c@{}}
\toprule
 & & \multicolumn{2}{c}{$H$} & \multicolumn{2}{c}{$H_\mathrm{loc}(1000~\mathrm{s})$} & & & & & \\
\cmidrule(lr){3-4}\cmidrule(lr){5-6}
pilots & flights & $\G_n$ & alone & $\G_n$ & alone & \makecell{time in\\climb (\%)} & \makecell{searches\\to climb (\%)} & \makecell{climb\\rate (m/s)} & \makecell{glide speed\\(km/h)} & \makecell{glide\\ratio} \\
\midrule
solo & 119\,129 & 0.899 &  & 0.837 &  & 30.8 & 43.3 & 1.22 & 30.1 & 10.4 \\
2 & 26\,411 & 0.917 & 0.914 & 0.889 & 0.861 & 32.7 & 49.0 & 1.27 & 33.2 & 12.5 \\
3 & 9\,480 & 0.923 & 0.919 & 0.906 & 0.867 & 33.6 & 49.9 & 1.26 & 33.5 & 13.1 \\
4 & 4\,073 & 0.924 & 0.920 & 0.910 & 0.865 & 33.6 & 51.0 & 1.25 & 33.5 & 13.7 \\
5--6 & 2\,686 & 0.925 & 0.922 & 0.915 & 0.870 & 33.1 & 50.1 & 1.23 & 33.1 & 14.1 \\
7+ & 1\,142 & 0.933 & 0.924 & 0.944 & 0.871 & 34.6 & 51.7 & 1.20 & 34.6 & 15.9 \\
\bottomrule
\end{tabular}

\caption{Paragliders, pooled by the size of the configuration (``alone'': the same flights alone).}
\label{tab:size}
\end{table}

\begin{figure}
\centering
\includegraphics[width=\textwidth]{decompose_size.pdf}
\caption{Orange: $\G_n-\A$ within the flight. Blue: $\G_n-\G_2$ at the same site-day and hour. Values
with bands in Table~\ref{tab:size-contrasts}.}
\label{fig:size}
\end{figure}

Pooled, $\Hloc(1000~\mathrm{s})$ rises with $n$, from \HlocThousandTwo{} with 2 pilots to
\HlocThousandSeven{} with 7+, while the same flights alone stay between \HlocThousandAloneMin{} and
\HlocThousandAloneMax{} (Figure~\ref{fig:size-hloc}, Table~\ref{tab:size}). Within the flight the
effect also grows, from \DHlocThousandFlightTwo{} with 2 pilots to \DHlocThousandFlightSeven{} with
7+ (Figure~\ref{fig:size}); there the day and the site are fixed, so the growth comes from the hour or
from the flights that join large groups. At the same site-day and hour it vanishes up to six pilots:
configurations of 3 to 6 pilots differ from pairs by \DHlocThousandMiddleMin{} to
\DHlocThousandMiddleMax, with bands across zero. Only 7+ stands apart (\DHlocThousandSevenTwoSdh{}
\DHlocThousandSevenTwoSdhBand, cross-country speed \DXcSevenTwoSdh{} km/h), a rare configuration, 3\%
of the group time. These may be competition gaggles flying a set task; the archive does not flag
competitions, so this stays uncontrolled. Time in climb and the flatter glide do grow with $n$ at the
same place and hour, from \DClimbSizeMin{} to \DClimbSizeMax{} points and from \DSinkSizeMin{} to
\DSinkSizeMax{} m/km. On six single site-days the differences between days exceed those between
sizes (Appendix~\ref{app:sitedays}).

\subsection{Leader and follower}\label{sec:mechanisms}

In configurations of two pilots, the signed median delay $\ell_{\rm med}$ of the pair interval
(Appendix~B of the note) says who passes first: the leader, when the other passes at least 15~s later;
otherwise the two fly abreast (\Abreast\% of the pairs). Leader and follower are compared within the
same configuration, the configuration being the stratum of Eq.~\eqref{eq:contrast}
(Appendix~\ref{app:leaders}). The pilot ahead climbs more and a little faster and glides faster (time
in climb \LeadClimb{} points, climb rate \LeadRate{} m/s, glide speed \LeadSpeed{} km/h), and glides
as flat (\LeadSink{} m/km). The share of searches that end in a climb does not differ (\LeadFound{}
points): on this measure the follower shows no sign of finding thermals through the leader. How
groups form and break and whether the same pilots lead are in section 7 of the note.

\paragraph{What we keep.} Most of the effect is there with two pilots. Beyond a pair, at the same
site-day and hour, size adds climbing and flatter glides but no straighter progress, except in the
rare configurations of 7 or more. The leader is a little faster, and ends searches in a climb as often
as the follower.
